\scanpage{60}
est $I$, et on a un \deleted{\unsure{mono}} \uline{morphisme}
\[
 I\hookrightarrow X \tag*{(26)}
\]
\deleted{stable \unsure{avec|sous} l'acti} compatible avec l'action de $\underline V$, d'où on déduit facilement que c'est un \uline{monomorphisme} (car on vérifie que $I$ s'envoie dans l'ouvert \deleted{\unsure{o}} $\underset{(27)}{X^*}$ de $X$ où l'action de $\underline V$ est libre, cf plus bas).

Ainsi on est dans la situation du § 2, qui implique que $\Sigma\simeq\Sigma_{\underline J}$, où $\underline J=\underline V^*$, que $I$ est l'image inverse d'une section bien déterminée $t$ de $\Sigma^*_{\underline J}=\Sigma_{\underline J}\smallsetminus D_{\underline J}$, l'image inverse $\Delta$ de $\underline J$ dans $X$ étant \deleted{\unsure{une}} la configuration \unsure{icosaédrale|icosaèdrale} dans $X$, canoniquement associée à l'action donnée de $\underline V$ sur \deleted{le} fibré en droites projectives $\underline V$ sur $X$.

En termes de l'action de $\underline{\deleted{\Delta}V}$ sur $X$, $\underline\Delta$ se récupère comme le diviseur de ramification pour cette action -- ensemblistement, il correspond donc aux pts (à val. dans des corps alg. clos) dont le stabilisateur ds $V$ est $\ne 0$, i.e. il correspond à des pts $x\in E$ tels qu'il existe $\unsure{\zeta|\xi}\in V={}_{\unsure{2|8}}E_0$, $\zeta\ne0$, avec $x+\zeta\in\{x,\sigma x=a-x\}$ i.e $x+\zeta=a-x$ i.e. $2x+\zeta=a$ d'où $4x=2a$ \deleted{i.e. $x\in \unsure{I\mathbin{\wedge}{}_4E_0}$,} \inserted{mais $2x\ne a$}
% archive sheet number: 59
